\section{结算机制与采购行为}
\label{sec:settlement_mechanism}

\subsection{先固定结算口径，再比较策略}

赛题规定计划购电按计划量计费，上调部分按交易电价的1.5倍支付，下调部分按交易电价的50\%承担违约费用，紧急购电按5倍电价支付。对取消的原计划是否返还原购电费，文字没有另行给出退款项。本文主方案采用原计划全额付费且不退款的解释，同时将退款解释作为敏感性对照。两种解释下的单时段非紧急费用见表\ref{tab:mech_settlement_convention}，其中$[x]_+=\max\{x,0\}$，$g$为零时原计划，$a$为最终有效常规购电量。

\begin{table}[htbp]
\centering
\caption{结算口径对比与单时段示例（不含紧急购电费）}
\label{tab:mech_settlement_convention}
\small
\begin{tabularx}{\linewidth}{p{0.18\linewidth}Xrr}
\toprule
口径 & 非紧急费用 & 下调示例 & 上调示例\\
\midrule
主口径：不退款 & $pg+1.5p[a-g]_++0.5p[g-a]_+$ & 110元 & 130元\\
退款敏感性口径 & $p\min\{g,a\}+1.5p[a-g]_++0.5p[g-a]_+$ & 90元 & 130元\\
\bottomrule
\end{tabularx}
\par\smallskip
\begin{minipage}{\linewidth}
\footnotesize 注：两例均取$g=100$ kWh、$p=1$元/kWh；下调例取$a=80$ kWh，上调例取$a=120$ kWh。退款解释仅在下调时返还取消部分的原计划费用。
\end{minipage}
\end{table}

在任一口径下，紧急费用均另加$5ph$，多次修订均与零时原计划比较，并按最后生效的购电量结算，不对相邻修订重复收费。历史负载、光伏和价格数据只能用于评估给定结算定义下的策略结果，不能凭数值结果反证结算定义，更不能因退款解释费用较低就选择该口径。后文的主结果与退款敏感性结果分别汇报。

相同修订策略在两种口径下的334天费用相差320347.51元，其量级大于若干控制增强的收益。因此先固定结算语义，再解释模型改进；完整敏感性结果见表\ref{tab:v2_settlement}。

\paragraph{不退款口径下的下调支配性。}
假设$p>0$、原计划不退款且富余电量允许免费弃置。若某可行解有$a<g$，把有效购电量改为$a'=g$，保持原充放电、紧急购电及储电轨迹不变，将新增的$g-a$电量弃置，物理约束仍满足，结算费用减少$0.5p(g-a)>0$。因此至少存在一个满足$a\ge g$的最优解。该结论依赖上述三个条件，在表\ref{tab:mech_settlement_convention}的退款解释下不能直接使用。对固定购电量的内层储能控制器而言，此结论也不授权其修改外层已经承诺的购电量。

\subsection{五倍补购价格导出的80\%采购分位数}

为解释采购的经济动机，先剥离储能及跨时段约束。设$X$为单时段实际净需求电量，可以为负；$q\ge0$为提前承诺的采购量，$p>0$为确定电价，富余电量无残值，且$\mathbb E[X_+]<\infty$。给定决策时已有的信息，以下概率和期望均针对当时可得的条件分布。期望结算费用为
\begin{equation}
 J_2(q)=pq+5p\,\mathbb E[(X-q)_+].
 \label{eq:mech_q2_cost}
\end{equation}
记$F_X(x)=\Pr(X\le x)$、$F_X(x^-)=\Pr(X<x)$。函数$J_2$为凸函数，其次梯度区间为
\begin{equation}
 \partial J_2(q)=p[5F_X(q^-)-4,\ 5F_X(q)-4].
 \label{eq:mech_q2_subgradient}
\end{equation}
对于$q>0$的内点最优解，充要条件是
\begin{equation}
 F_X(q^-)\le0.8\le F_X(q).
 \label{eq:mech_q2_quantile_condition}
\end{equation}
这同时适用于连续分布和含原子的离散场景分布。若分布在最优点连续，则有$J_2'(q)=p-5p\Pr(X>q)=0$，即$F_X(q)=0.8$；若最优点位于$q=0$边界，则右导数非负的条件为$F_X(0)\ge0.8$。定义下分位数$F_X^{-1}(\alpha)=\inf\{x:F_X(x)\ge\alpha\}$，可选取
\begin{equation}
 q_2^\star=\max\{0,F_X^{-1}(0.8)\}.
 \label{eq:mech_q2_quantile}
\end{equation}
若分布存在平坦区间，最优解可能不唯一，该式给出其中一个可行选择。

因此，五倍补购价格本身就使低估需求的经济损失高于高估需求的损失。即便仅最小化期望费用、完全不加入CVaR，最优采购也可能高于均值。采购偏高不能单独作为风险厌恶的证据，风险扩展是否值得保留仍应由额外费用及其带来的可观察收益判断。

\subsection{1.5倍增购价格导出的70\%条件阈值}

在主结算口径下，固定零时计划$g\ge0$，由下调支配性只需研究$a\ge g$的增购决策。将已经承诺的$pg$视为常数，当前剩余费用为
\begin{equation}
 J_3(a)=1.5p(a-g)+5p\,\mathbb E[(X-a)_+],\qquad a\ge g.
 \label{eq:mech_q3_cost}
\end{equation}
同理，$\partial J_3(a)=p[5F_X(a^-)-3.5,\ 5F_X(a)-3.5]$。当$a>g$时，最优条件为$F_X(a^-)\le0.7\le F_X(a)$；在$a=g$边界，条件为$F_X(g)\ge0.7$。因此可取
\begin{equation}
 a_3^\star=\max\{g,F_X^{-1}(0.7)\}.
 \label{eq:mech_q3_quantile}
\end{equation}
式\eqref{eq:mech_q3_quantile}中的分布必须由修订时已经获得的信息确定。70\%是给定原计划后的单时段增购阈值，不是问题三零时计划的最优分位数，也不表示增购量直接等于需求的70\%分位数。实际增购量为$[F_X^{-1}(0.7)-g]_+$。零时计划如要计入以后修订的价值，还需描述未来信息到达及合法修订决策，不能由这一单时段推导替代。

\subsection{波动价格下的价格加权需求分位数}

若未来结算电价$P$也未知，设$P>0$几乎处处成立，$0<\mathbb E[P]<\infty$，且$\mathbb E[PX_+]<\infty$。单时段日前决策在当时信息下固定$q$，期望费用为
\begin{equation}
 J_4(q)=\mathbb E[P]q+5\mathbb E[P(X-q)_+],\qquad q\ge0.
 \label{eq:mech_q4_cost}
\end{equation}
定义价格加权分布
\begin{equation}
 G(x)=\frac{\mathbb E[P\,\mathbf 1_{\{X\le x\}}]}{\mathbb E[P]},\qquad
 G(x^-)=\frac{\mathbb E[P\,\mathbf 1_{\{X<x\}}]}{\mathbb E[P]}.
 \label{eq:mech_weighted_cdf}
\end{equation}
由$P>0$和有限期望，$G$是合法分布函数，并有
\begin{equation}
 \partial J_4(q)=\mathbb E[P][5G(q^-)-4,\ 5G(q)-4].
 \label{eq:mech_q4_subgradient}
\end{equation}
内点最优条件为$G(q^-)\le0.8\le G(q)$，在连续点化为$\mathbb E[P\mathbf 1_{\{X\le q\}}]/\mathbb E[P]=0.8$。考虑非负约束，可取
\begin{equation}
 q_4^\star=\max\{0,G^{-1}(0.8)\}.
 \label{eq:mech_weighted_quantile}
\end{equation}
在离散场景$(x_s,p_s,\pi_s)$下，$p_s>0$、$\pi_s\ge0$、$\sum_s\pi_s=1$，只需把场景概率改写为
\begin{equation}
 \widetilde\pi_s=\frac{\pi_sp_s}{\sum_j\pi_jp_j},\qquad
 G(x)=\sum_{s:x_s\le x}\widetilde\pi_s.
 \label{eq:mech_scenario_reweight}
\end{equation}
若$P$与$X$独立，则$G=F_X$，退化为普通80\%分位数；若高需求经常伴随高价格，相应场景在采购决策中的权重增大。该机制说明，问题四应保留价格与净负荷的联合关系，分别打乱两类历史残差会改变优化问题所使用的经济权重。

\subsection{与储能联合优化及预测评价的衔接}

以上阈值是无储能单时段子问题的解析基准。加入储能后，当前多购或少购的影响会通过充放电效率、容量、功率及后续电价传递到其他时段，电池存量具有随时间和状态变化的机会价值。因此，80\%和70\%不能直接逐时套用于完整调度，也不能为参考备用公式中的80\%误差分位数提供理论背书。参考备用参数和采购分位数具有不同的决策对象与约束作用，应分别检验。

本文据此将模型递进关系明确为：结算规则决定非对称误差成本，场景采购刻画这种成本，再由储能联合优化处理跨时段耦合，实际控制在已经承诺的常规购电量和合法信息范围内执行。经济主方案以实际总购电费用为评价依据，CVaR和参考备用作为可选运行偏好，与期望费用策略在相同条件下比较。终端价值近似系数记为$\beta_T$，单位为元/kWh；备用缩放系数记为$k_\rho$，为无量纲数，二者不复用符号。

预测评价也应对应这一经济目标。文献\cite{elmachtoub2017spo}提出以预测引致的决策损失评价预测，并进一步构造SPO与SPO+方法。
% Citation locator: https://arxiv.org/abs/1710.08005, Abstract paragraphs 1--2.
本文仅借鉴其决策导向思想：在相同历史信息、初始储电量、场景构造和控制器下，将候选预测器输入统一调度程序，以一月逐日实际结算费用选择预测方法，RMSE和MAE作为描述性指标。在预测器选型后再比较控制器，属于有限候选集上的分步验证，不等同于实现SPO或SPO+训练算法。2月至12月回测费用用于评价冻结后的规则，不能据此事后挑选最低费用方案；修订研究已经参考过旧版该时段结果，也不把它称为从未参与研究过程的独立测试。
